\documentclass[11pt]{report}
\usepackage{geometry}
\usepackage{graphicx}
\usepackage{titling}
\usepackage{fancyhdr}
\usepackage{enumerate}
\usepackage{ulem}
\usepackage{latexsym}
\usepackage{float}
\usepackage{amssymb}
\usepackage{amsthm}
\usepackage{amsfonts}
\usepackage{amsmath}
\usepackage[labelfont=bf]{caption}
\usepackage[usenames,dvipsnames,svgnames,table]{xcolor}
\usepackage{listings}
\usepackage{lmodern}
\usepackage{../Maple/maplestd2e}
\parindent=0pt
%\frenchspacing

\pagestyle{fancy}

\fancyhead[L]{\slshape\footnotesize April 7, 2014\\\textsc{01410 Cryptology 1}}
\fancyhead[R]{\slshape\footnotesize \textsc{Andreas Kjeldsen (s092638)}\\\textsc{Morten Eskesen (s133304)}}
\fancyfoot[C]{\thepage}
\newcommand{\HRule}{\rule{\linewidth}{0.5mm}}

\definecolor{lightgrey}{rgb}{0.975,0.975,0.975}
\lstdefinestyle{maple}{
	morekeywords={and,assuming,break,by,catch,description,do,done,
	elif,else,end,error,export,fi,finally,for,from,global,if,
	implies,in,intersect,local,minus,mod,module,next,not,od,
	option,options,or,proc,quit,read,return,save,stop,subset,then,
	to,try,union,use,uses,while,xor},
	sensitive=true,
	morecomment=[l]\#,
	morestring=[b]",
	morestring=[d]",
	frame=single,
 	basicstyle=\tt\small,
	keywordstyle=\bf,
	xleftmargin=.25in,
	xrightmargin=.25in,
	mathescape=true,
	backgroundcolor=\color{lightgrey},
}[keywords,comments,strings]

\begin{document}

\begin{titlepage}
\begin{center}

\includegraphics[scale=2.0]{../GFX/dtu_logo.pdf}\\[1cm]
\textsc{\LARGE Technical University of Denmark}\\[1.5cm]
\textsc{\Large 01410 Cryptology 1}\\[0.5cm]

% Title
\HRule \\[0.4cm]
{\huge \bfseries Homework 2}\\[0.1cm]
\HRule \\[1.5cm]

% Author and supervisor
{\large
\emph{Authors:} \\[10pt]
Andreas Hallberg \textsc{Kjeldsen}\\
\emph{s092638@student.dtu.dk} \\[10pt]
Morten Chabert \textsc{Eskesen}\\
\emph{s133304@student.dtu.dk}
}
\vfill

% Bottom of the page
{\large April 7, 2014}

\end{center}
\end{titlepage}

\section*{Exercise 2.1}
\subsection*{2.1.1}
We have to show that $m^{e\tilde{d}} \equiv m$ mod $n$ for all $m \in \mathbb{Z}_{n}$.\\
The keys $e$ and $\tilde{d}$ are chosen such that.
$$e\tilde{d} \equiv 1 \text{ mod } \frac{(p-1)(q-1)}{gcd(p-1,q-1)}$$

This means that for some positive integer $k$
$$e\tilde{d} = 1 + k\frac{(p-1)(q-1)}{gcd(p-1,q-1)}$$

We can rewrite this expression and get
$$m^{e\tilde{d}} \text{ mod } n = m^{1+k\frac{(p-1)(q-1)}{gcd(p-1,q-1)}}  \text{ mod } n$$
for some integer $k$.\\
\\
If two integers $x$ and $y$ are congruent modulo $n$ then they are also congruent modulo $p$ and modulo $q$ because both $p$ and $q$ divide $n$. The Chinese Remainder Theorem tells us that the reverse implication is also true. This means that if $x$ and $y$ are congruent modulo $p$ and congruent modulo $q$, then they are also congruent modulo $n$.\\
\\
We want to show that $m^{e\tilde{d}} \equiv m$ mod $n$ so it will be sufficient to show that:

$$m^{e\tilde{d}} \equiv m \text{ mod } p \text{ and } m^{e\tilde{d}} \equiv m \text{ mod } q$$

First we will show that $m^{e\tilde{d}} \equiv m$ mod $p$. We therefore have two cases to consider:
\begin{enumerate}
	\item $p$ divides m
	\item $p$ does not divide m.
\end{enumerate}

\begin{description}
	\item[Case 1:] If $p$ divides $m$, then $m \equiv 0$ mod $p$, but also $m^{e\tilde{d}} \equiv 0$ mod $p$, therefore $m^{e\tilde{d}} \equiv m$ mod $p$.
	
	\item[Case 2:] If $p$ does not divide m then $m \in \mathbb{Z}^*_p$. By Fermat's Little Theorem we have $m^{p-1} \equiv 1$ mod $p$. Since $e\tilde{d} \equiv 1$ mod $\psi(n)$, we have that $\psi(n)$ divides $e\tilde{d} - 1$.\\
	This gives:
	\begin{center}
		$k \psi(n) = e\tilde{d} - 1$, so $e\tilde{d} = k \psi(n) + 1$ for some integer $k$.
	\end{center}
	We therefore have:
	$$m^{e\tilde{d}} = m^{k\psi(n)+1} = m * m^{k\frac{(p-1)(q-1)}{gcd(p-1,q-1)}}$$
	$$m^{e\tilde{d}} = m * (m^{p-1})^{k\frac{(q-1)}{gcd(p-1,q-1)}}$$
	$$m^{e\tilde{d}} \equiv m * 1^{k\frac{(q-1)}{gcd(p-1,q-1)}} \text{ mod } p$$
	$$m^{e\tilde{d}} \equiv m \text{ mod } p$$
\end{description}

We can do similar calculations to show that $m^{e\tilde{d}} \equiv m$ mod $q$ by replacing $p$ by $q$.\\
Therefore we have now shown for all $m \in \mathbb{Z}_n$ that
$$m^{e\tilde{d}} \equiv m \text{ and } p \text{ and } m^{e\tilde{d}} \equiv m \text{ mod } q$$

Hence we can concluce that $m^{e\tilde{d}} \equiv m$ mod $n$ for all $m \in \mathbb{Z}_{n}$.


\subsection*{2.1.2}
Let $p = 881$, $q = 461$, and let $n = pq = 405141$. We have to show that $e = 3$ is an allowed encryption exponent for an RSA encryption system with modulus $n$. By the definition of RSA $e$ must be chosen such that $e$ and $\phi(n)$ are co-prime. Formally this means that $e \in \mathbb{Z}^*_{\phi(n)}$, where $\phi(n) = (p-1)(q-1)$.
$$gcd(3, (881-1)(461-1) = 1$$
This means that $e$ and $\phi(n)$ are co-prime and therefore $e=3$ is an allowed encryption exponent.

\subsection*{2.1.3}
We have to find $d_1$ such that $ed_1 \equiv 1$ mod $\phi(n)$.
\begin{mapleinput}
p := 881; q := 461; e := 3; \\
d1 := mod($e^{-1}$, (p-1)*(q-1))
\end{mapleinput}

\begin{lstlisting}[style=Maple]
> p := 881; q := 461; e := 3;
d1 := mod($e^{-1}$, (p-1)*(q-1))
\end{lstlisting}

Using the maple code above we find that $d_1 = 269867$
\subsection*{2.1.4}
We have to find $d_2$ such that $ed_2 \equiv 1$ mod $\psi(n)$

\begin{mapleinput}
p := 881; q := 461; e := 3; \\
d2 := mod($e^{-1}$, (p-1)*(q-1)/gcd(p-1, q-1))
\end{mapleinput}

\begin{lstlisting}[style=Maple]
> p := 881; q := 461; e := 3;
d2 := mod($e^{-1}$, (p-1)*(q-1)/gcd(p-1, q-1))
\end{lstlisting}

Using the maple code above we find that $d_2 = 6747$.
\subsection*{2.1.5}
Choosing $\psi(n)$ instead of $\phi(n)$ in the congruence for d means that the decryption becomes faster since
$$lcm(p-1,q-1)=\frac{(p-1)(q-1)}{gcd(p-1,q-1)} \leq \frac{(p-1)(q-1)}{2}$$
Because $p$ and $q$ are odd primes $gcd(p-1,q-1) \geq 2$. 

\section*{Exercise 2.2}
\subsection*{2.2.a}
We have implemented trial division in maple with the following code:

\begin{mapleinput}
TrialDivision := proc (n::integer) \\
local i; \\
if n $\leq$ 1 then false \\
elif n = 2 then true \\
elif type(n, 'even') then false \\
else for i from 3 by 2 while i*i $\leq$ n do \\
if irem(n, i) = 0 then return false end if \\
end do; \\
true end if \\
end proc:
\end{mapleinput}

\begin{lstlisting}[style=Maple]
> TrialDivision := proc (n::integer)
local i;
if n $\leq$ 1 then false
elif n = 2 then true
elif type(n, 'even') then false
else for i from 3 by 2 while i*i $\leq$ n do
if irem(n, i) = 0 then return false end if
end do;
true end if
end proc:
\end{lstlisting}

\begin{mapleinput}
result := 0; \\
for n from 25 to 25000 do \\
if TrialDivision(n) then result := result+1 end if \\
end do; \\
result;
\end{mapleinput}

\begin{lstlisting}[style=Maple]
> result := 0;
for n from 25 to 25000 do
if TrialDivision(n) then result := result+1 end if
end do;
result;
\end{lstlisting}

Using this code we find that the number of primes $s$ between 25 and 25000 is 2753.
$$s = 2753$$

\subsection*{2.2.b}
We have implemented the Miller-Rabin algorithm with $k$ iterations in maple with the following code:
\begin{mapleinput}
MillerRabin := proc (n::integer, k::integer) \\
local x, r, roll, s, d, i, a; \\
s := n-1; d := 0; \\
while mod(s, 2) = 0 do \\
s := (1/2)*s; d := d+1 \\
end do; \\
for i to k do \\
roll := rand(2 .. n-1); \\
a := roll(); x := mod($a^s$, n); \\
if x = 1 or x = n-1 then next end if; \\
for r to d-1 do x := mod($x^2$, n); \\
if x = 1 then return false end if; \\
if x = n-1 then break end if \\
end do; \\
if x $\ne$ n-1 then return false end if \\
end do; \\
return true \\
end proc:
\end{mapleinput}

\begin{lstlisting}[style=Maple]
> MillerRabin := proc (n::integer, k::integer)
local x, r, roll, s, d, i, a;
s := n-1; d := 0;
while mod(s, 2) = 0 do
s := (1/2)*s; d := d+1
end do;
for i to k do
roll := rand(2 .. n-1);
a := roll(); x := mod($a^s$, n);
if x = 1 or x = n-1 then next end if;
for r to d-1 do x := mod($x^2$, n);
if x = 1 then return false end if;
if x = n-1 then break end if
end do;
if x $\ne$ n-1 then return false end if
end do;
return true
end proc:
\end{lstlisting}

\subsection*{2.2.c}
We use this maple code below and define $k = 1,2,3,\dots$ to find the smallest number of iterations needed such that we gets the correct answer $s$.
\begin{mapleinput}
result := 0; \\
for n from 25 to 25000 do \\
if MillerRabin2(n, k) then result := result+1 end if \\
end do; \\
result;
\end{mapleinput}

\begin{lstlisting}[style=Maple]
> result := 0;
for n from 25 to 25000 do
if MillerRabin2(n, k) then result := result+1 end if
end do;
result;
\end{lstlisting}

This gives us the following table:
\begin{center}
\begin{tabular}{| c | c | c | c | c |}
\hline
k & 1 & 2 & 3 & 4 \\ \hline
s & 2792 & 2755 & 2754 & 2753 \\ \hline
\end{tabular}
\end{center}
With $k = 4$ iterations we get the correct answer for $s$ which is 2753.
\end{document}
